\section{Notation from the source-cover paper}\label{sec:preliminaries}

We use the notation of~\cite{SourceCovers2026} and recall only what the statements
below need. All objects in this section are defined and Lean-verified there.

\paragraph{Geometries.}
Fix integers $d,b\ge1$ and put $L=d^2b$, $W=db$ and $\Omega=\{0,\dots,L-1\}$; for~(T),
$b=b(d)$. A fixed complete geometry consists of the input-head and work-head walks,
the halting time $T$ and the append times $t_0<\dots<t_{L-1}<T$
\cite[Section~2.1]{SourceCovers2026}. The \emph{raw path} $x$ is the input-head
position clamped to $\Omega$; it moves by at most one cell per step. In the Lean model
times are natural numbers, and statements about all times quantify over all of them.
The first append time of column $j$ is $s_j=t_{jW}$
\cite[Theorem~2.9]{SourceCovers2026}. An \emph{admissible geometry} with parameters
$d,b,c,\beta,N,q,a_0$ \cite[Definition~2.2]{SourceCovers2026} adds a nonempty
correctness set $\mathcal C$, the runtime bound $T\le cL$, whole-step record capacity
at most $\beta$, an $N$-class span partition of the full family system, an original
$\mathrm Q_2$ core $\mathcal K_0$, a disjoint set $\mathcal E$ of at most $q$
exceptional complete families, and the bound $a_0$ on the hop count of every other
family. Its correctness deficit is $\Delta=L-\log_2|\mathcal C|$
\cite[(2.2)]{SourceCovers2026}.

\paragraph{Families and shadows.}
A complete actual family $F$ has an owner $o_F$, an outermost hop $(p_F,U_F)$, a full
span $(p_F,U_F]$ and a hop count $|F|$ \cite[Definition~2.1]{SourceCovers2026}.
We call $(p_F,o_F]$ its \emph{lower region} and $(o_F,U_F]$ its \emph{active
interval}. Its clipped shadow $s_F$ is the identity off the active interval and moves
every time of it; by \cite[Definition~2.4 and (2.6)]{SourceCovers2026},
\begin{equation}\label{eq:strict-shadow}
 p_F<s_F(t)\le o_F<t\le U_F\qquad\text{whenever }o_F<t\le U_F;
\end{equation}
on its active interval $s_F$ is antitone. We write $F(t)=s_F(t)$ and say that $F$
\emph{acts} at $t$ if $F(t)\ne t$, that is, if $t$ lies in its active interval.
For a set $\mathcal S$ of families, $\mathcal O_{\mathcal S}(t)$ is the set of values
obtained from $t$ by finite compositions of their shadows, including the empty
composition \cite[Definition~2.4]{SourceCovers2026}; its elements are at most $t$.
A word $(F_1,\dots,F_n)$ of families is \emph{successful} from $t$ if $F_1$ acts at
$t$, $F_2$ acts at $F_1(t)$, and so on; its \emph{endpoint} is
$F_n(\cdots F_1(t)\cdots)$.

\paragraph{Cores, phases and banks.}
A set $\mathcal K$ of complete families is $\mathrm Q_2$ if its spans split into two
classes of pairwise disjoint spans. Joining two families of $\mathcal K$ when their
spans intersect gives a caterpillar forest, the \emph{overlap forest}; its
\emph{original nonleaves} are the vertices with at least two neighbours. The nonleaves
of each component form a path, so a nonleaf has at most two nonleaf neighbours, and a
leaf has at most one neighbour \cite[Lemma~2.6]{SourceCovers2026}. Fix $M\ge2$ and a
protected set of at most $R<M$ families of $\mathcal K$. For a residue
$\rho\in\{0,\dots,M-1\}$, the original nonleaves whose rank along their nonleaf path
is congruent to $\rho$ modulo $M$ are deleted, and the other families of $\mathcal K$
\emph{survive}; they form the retained core. A residue is a \emph{legal phase} if no
protected nonleaf has a rank congruent to it; there are at least $M-R$ legal phases.
Let $A_\rho(u)$ be the orbit of the retained core at $u$. Every successful word of
surviving families has at most $M+1$ letters, and $|A_\rho(u)|\le2M+3$
\cite[Lemma~2.6]{SourceCovers2026}. At most two families of $\mathcal K$ act at any one
time, since their spans contain that time and so lie in distinct classes. For an
exceptional set $\mathcal E$, the exact source bank has value set $O_\rho(t)$ at $t$,
the orbit of the retained core together with the exceptional families
\cite[Theorem~2.7]{SourceCovers2026}. The raw source set is $S_\rho(t)=x(O_\rho(t))$
\cite[Section~3.1]{SourceCovers2026}; it contains $x(t)$. In limit statements,
$\mathcal K$ consists of the first $J$ families of the original core in its supplied
order by nonincreasing hop count, the protected set consists of the first $R$, and the
forest, ranks and phases are rebuilt for every selection
\cite[Section~4.2]{SourceCovers2026}.

\paragraph{Gap area and the paid error.}
View $\Omega$ as the integer points of a circle of circumference $L$. For a nonempty
$S\subseteq\Omega$ whose consecutive distinct points have circular gaps
$g_1,\dots,g_v$, the point-gap area is \cite[(2.14)]{SourceCovers2026}
\[
 \gaparea(S)=\frac1{2L^2}\sum_{h=1}^vg_h^2\in\Bigl[0,\frac12\Bigr].
\]
We call $\gaparea(S_\rho(s_j))$ the \emph{frozen point area} of phase $\rho$ and
column $j$. One counting fact is used repeatedly. If a nonempty $S\subseteq\Omega$
lies in the union of at most $n$ closed integer intervals $[a_i,b_i]$ with
$\sum_i(b_i-a_i)\le\alpha L$, then
\begin{equation}\label{eq:interval-budget}
 \gaparea(S)\ge\frac{(1-\alpha)_+^2}{2n},\qquad (z)_+=\max(z,0).
\end{equation}
This is the exposed-gap step in the proof of \cite[Lemma~3.3]{SourceCovers2026}:
an interval contains the left endpoint of at most one gap that it does not contain,
the exposed gaps therefore have total length at least $(1-\alpha)_+L$, and the
Cauchy--Schwarz inequality applies to them. Its Lean form is
\lean{LeanMachine.IntervalGapArea.point_area_lower_of_budgets} in the shared library.
The paid error of \cite[(3.7)]{SourceCovers2026} is
\[
 \mathsf E_q(w;J,M,R,a,\tau)=\frac{2C_q(T-1)}{dL}+\frac{2(w-1)}L+\frac{W-b+1}{2L}
 +\frac\beta L\Bigl[J+q+\frac\tau{M-R}+\frac{(1+Na)C_q(T-1)}w\Bigr],
\]
with $C_q=C_0(1+3C_0)^q$ and $C_0=16+24M(M-1)$.

\paragraph{Near-full sequences and joint profiles.}
A near-full sequence consists of admissible geometries indexed by $\nu$, with common
$c,\beta,N,q$, such that $d_\nu\to\infty$, $a_{0,\nu}/L_\nu\to0$ and
$\Delta_\nu/L_\nu\to0$ \cite[Definition~4.1]{SourceCovers2026}. Each geometry comes
with an order $G_1,G_2,\dots$ of its original core by nonincreasing hop count.
A joint profile is a coordinatewise limit $(h_i,v_i)_{i\ge1}$ of the normalized
profiles $h_i^{(\nu)}=|G_i|/L_\nu$ and $v_i^{(\nu)}=\Var(x_\nu;[p_{G_i},U_{G_i}])/L_\nu$
along one strictly increasing subsequence; both profiles are zero after the last
family. Write
\[
 \tau_R=\sum_{i>R}h_i,\qquad \tau_{R,J}=\sum_{R<i\le J}h_i,\qquad V_K=\sum_{i>K}v_i.
\]
The limiting profiles are summable, so $\tau_R\to0$ as $R\to\infty$. Limits in $\nu$
are taken along the subsequence of the profile. A strictly increasing subsequence of a
near-full sequence is near-full, and every near-full sequence has a joint profile
\cite[Section~4.1]{SourceCovers2026}.
